\subsection{\label{subsecunidentified}Unidentified Charged Particles}



Spectra of positively and negatively charged particles were obtained separately and summed afterwards.
The differences between the final spectra for particles and antiparticles were found to be around ~1.5\%.
The unidentified charged particles were reconstructed using the combined information from ITS and TPC. The \pt~range of the spectra in all multiplicity classes based on the V0M amplitudes is 0.16-40\ GeV/$c$ and the pseudorapidity was limited to $|\eta| < 0.5$. 
This pseudorapidity limit allows a comparison of the charged hadron spectrum with the sum of pions, kaons and protons analyzed in the rapidity range $|y| < 0.5$ by transforming them to the corresponding pseudorapidity window with the appropriate Jacobian ${{\rm d}^2{N} \over {\rm d}y{\rm d}\pt} = {E \over p} {{\rm d}^2{N} \over {\rm d}\eta{\rm d}\pt}$ for each \pt-interval. This cross-check showed a difference of less than 5\% which is consistent within the non-common systematics with the expected contributions from electrons, muons, and heavier charged baryons that are counted in addition to pions, kaons, and protons in the charged hadron spectrum.
% see comparison plot in: https://alice-notes.web.cern.ch/system/files/notes/analysis/477/2018-Apr-27-analysis_note-Internal_Note_ch_part_20161214_v4.pdf (fig 36)

The contribution from secondary particles was calculated in the same manner as described in detail in Section \ref{subsecPiKP}. The additional corrections based on FLUKA \cite{Ferrari:2005zk,BOHLEN2014211} for kaons and anti-protons, which are needed for those specific identified particle measurements in order to account for an imperfect description of absorption cross-section in the detector material, were found to have negligible impact on the unidentified charged particle spectra and were therefore not applied. 

The systematic uncertainties are summarized in Table \ref{tab:sysunidentified}. The multiplicity dependence of the tracking efficiency and the feed-down correction were found to be less than 1\% and were included in the final systematic uncertainty. The total systematic uncertainty is \pt~dependent, with values around 6-7\% up to 20 GeV/$c$. It reaches 9.6\% for the highest \pt~bin. The main contributions to the total systematic uncertainty come from the global tracking efficiency (5\%) and the parameter variation for the track selection criteria (3-7\%). The other sources have a \pt~dependent contribution of less than 2\% each. The systematic uncertainties related to the dependence of the reconstruction efficiency on the MC event generator and the particle composition have been studied as described in \cite{Abelev:2013ala}. All sources of uncertainty are assumed to be uncorrelated and the total uncertainty was calculated as the quadratic sum of the different contributions. The systematic uncertainty contribution that is uncorrelated across multiplicities was 
estimated to be 2.1\% for all the V0M multiplicity bins over the entire \pt~range.


\begin{table*}[ph!]
  \begin{tabularx}{\textwidth}{p{5.3cm}*{3}{Y}}
\hline\hline
    {\bf Source} &  \multicolumn{3}{c}{Uncertainty (\%)}   \\
\hline
    \pt\ (\GeVc)   &   0.16 & 3.0 & 40.0  \\
    \cmidrule(l{5pt}){2-4} 

    Correction for secondaries & 0.2  & 0.2 & 1 \\
    Particle composition in secondaries & 1.7 & 1.3 & 0.8 \\
    Material budget & 1.5 & negl. & negl.  \\
    Global tracking efficiency & 5 & 5 & 5  \\
    Particle composition & 1.7 & 2 & 2 \\
    Track selection & 4 & 2.8 & 7.4 \\
    MC event generator & 1.1 & 1.8 & 2 \\
    \pt\ resolution & negl. & negl. & 0.5 \\
    Efficiency multiplicity dependence & 1 & 1 & 1 \\
\hline

    Total & 7.2 & 6.6 & 9.6  \\ 
    
    Total ($N_{ch}$-independent) & 6.9 & 6.3 & 9.4 \\ 
\hline\hline
 \end{tabularx}
 ~\newline
   \caption{Main sources and values of the relative systematic uncertainties of the \pt-differential yields for unidentified charged particles. The values are reported for low, intermediate and high \pt. The contributions that act differently in the various event classes are removed from the Total (quadratic sum of all contributions), defining the $N_{ch}$-independent ones, which are correlated across different multiplicity intervals.}
  \label{tab:sysunidentified}  
\end{table*}


